\section{Collective trimming of a viable subspace}
\label{sec:alvv_collective_trimming}

Arad, Landau, Vazirani and Vidick reduce the Schmidt ranks of a viable
subspace by trimming a basis of it collectively, cut by
cut~\cite[lines 642--697]{Arad2016RigorousRG}. Let $S$ be a subspace of
vectors on the first $i$ sites of a chain of $L$ sites and $T$ a subspace of
vectors on the chain. Write $\mathcal H_A$ for the first $i$ sites and
$\mathcal H_R$ for the remaining $L-i$, so that the cuts $(j:j+1)$,
$1\le j\le i-1$, of the source are the cuts of the chain after $j$ sites.

\begin{definition}[Trimmed subspace]
    \label{def:alvv_collective_trim}
    \lean{TNLean.MPS.AreaLaw.blockCut,
        TNLean.MPS.AreaLaw.chainTrimVector,
        TNLean.MPS.AreaLaw.chainTrimSubspace,
        TNLean.MPS.AreaLaw.chainTrimSubspace_eq_span}
    \leanok
    \uses{def:lvv_viable_set}
    Let $u_1,\dots,u_s$ be an orthonormal basis of $S$ and $\xi>0$, and set
    $\psi=\sum_k u_k\otimes|k\rangle\in\mathcal H_A\otimes\mathbb C^s$. For
    $j=1,\dots,i-1$ let $P^j$ be the orthogonal projection onto the left
    Schmidt vectors, across the cut after $j$ sites, of
    $(P^1\otimes1)\cdots(P^{j-1}\otimes1)\psi$ with Schmidt coefficient at
    least $\xi$. The trimmed vectors are
    $(P^1\otimes1)\cdots(P^{i-1}\otimes1)u_k$, and the trimmed subspace
    $\operatorname{Trim}_\xi(S)$ is their span. This is
    Definition~def:trimming of~\cite[lines 642--656]{Arad2016RigorousRG}.
\end{definition}

\begin{theorem}[Viability in block coordinates]
    \label{thm:alvv_collective_trim_coordinates}
    \lean{TNLean.MPS.AreaLaw.blockSplit,
        TNLean.MPS.AreaLaw.blockSplit_apply,
        TNLean.MPS.AreaLaw.leftSupportSubmodule_eq_comap,
        TNLean.MPS.AreaLaw.isCloseTo_leftSupportSubmodule_iff,
        TNLean.MPS.AreaLaw.blockCut_mono,
        TNLean.MPS.AreaLaw.mem_blockCut,
        TNLean.MPS.AreaLaw.blockCutSubtypeEquiv,
        TNLean.MPS.AreaLaw.blockRestSubtypeEquiv,
        TNLean.MPS.AreaLaw.rank_cutEnvVector_blockSplit}
    \leanok
    \uses{def:alvv_collective_trim}
    Write a vector $\psi$ on the chain in the coordinates
    $\mathcal H_A\otimes\mathcal H_R$. Then $\operatorname{Span}(S)\otimes
        \mathcal H_R$ is $\delta$-close to $T$ exactly when $S$ is $\delta$-viable
    for $T$ in these coordinates. For $1\le j\le i$, the Schmidt rank of $\psi$
    across the first $j$ sites of $\mathcal H_A$, with $\mathcal H_R$ on the
    right, equals the rank of the coefficient matrix $M_j(\psi)$ of the chain
    cut after $j$ sites.
\end{theorem}

\begin{proof}
    \leanok
    The coordinates are a relabelling of the configurations, so they preserve
    inner products, and $\psi$ lies in $\operatorname{Span}(S)\otimes
        \mathcal H_R$ exactly when every column of $M_i(\psi)$ lies in $S$. Both
    bipartitions in the second claim put the first $j$ sites on the left and
    the rest on the right, so the two coefficient matrices differ only by a
    relabelling of rows and columns.
\end{proof}

\begin{theorem}[Trimming a viable subspace on the chain]
    \label{thm:alvv_collective_trim}
    \lean{TNLean.MPS.AreaLaw.isCloseTo_leftSupportSubmodule_chainTrimSubspace,
        TNLean.MPS.AreaLaw.schmidtRank_cutVector_chainTrimVector_le}
    \leanok
    \uses{def:alvv_collective_trim, thm:alvv_collective_trim_coordinates}
    Let $S$ have dimension $s$ and let $\operatorname{Span}(S)\otimes
        \mathcal H_R$ be $\delta$-close to $T$. Suppose every vector of $T$ has
    $\operatorname{rank}M_j\le b$ at every cut $j=1,\dots,i-1$. Then for
    $\xi>0$ the subspace $\operatorname{Trim}_\xi(S)\otimes\mathcal H_R$ is
    $\big(\delta+2\xi\sqrt{(i-1)b}\big)$-close to $T$. Each trimmed vector
    has Schmidt rank at most $s/\xi^2$ across every cut $(j:j+1)$ of
    $\mathcal H_A$, $1\le j\le i-1$.
    This is Lemma~lem:trimming of~\cite[lines 660--697]{Arad2016RigorousRG},
    without the running time. The source states $\delta+\sqrt{ibs}\,\xi$, which
    its argument does not give; the bound above is the one it
    gives~\cite{gap:alvv16_trimming_constant}.
\end{theorem}

\begin{proof}
    \leanok
    In the coordinates of the previous theorem this is the trimming lemma for
    a subspace of a block of sites with the remaining sites as environment.
    The part of a vector of $\operatorname{Span}(S)\otimes\mathcal H_R$
    removed at the cut after $j$ sites pairs with a vector $t$ of $T$ to at
    most $\xi\sqrt b$ times the norm of the piece of $t$ removed there. The
    pieces removed at different cuts are orthogonal, so their norms sum to at
    most $\sqrt{i-1}\,\|t\|$. The rank bound holds because $\psi$ has at most
    $s/\xi^2$ Schmidt coefficients at least $\xi$ across each cut.
\end{proof}
